\section{Theoretical framework}
\label{sec:framework}
 
In this section, we review the necessary theoretical background
for the present calculation. We present the definitions of
CS kernel and rapidity evolution, and introduce the quasi-TMD wave functions. We then discuss the factorization of the quasi-TMDWFs
and its connection with the CS kernel. 
   
\subsection{Collins-Soper Kernel and Rapidity Evolution}


Unlike the collinear  lightcone  PDFs and  distribution amplitudes,  the TMDPDFs and TMDWFs depend on both the renormalization scale $\mu$ and an additional rapidity renormalization scale. The latter arises because the matrix elements also suffer from so-called rapidity divergences that require a dedicated regulator~\cite{Collins:1981uk,Becher:2010tm,Chiu:2011qc}.  In TMD factorizations, the contributions of hard, i.e. highly offshell, modes to the tree process are usually calculated in the dimensional regularization scheme. Collinear modes, which are related to  highly-boosted partons in distinct directions, and soft modes, whose typical momentum are  at the order  $\Lambda_{\mathrm{QCD}}$, share the same virtuality, %and are connected by Lorentz transformation. 
and are only distinguishable by their rapidity. In calculations using regularization schemes such as dimensional regularization, which only regulate ultra-violet divergences, one will encounter additional rapidity divergences that arise in soft and collinear matrix elements when integrating over rapidity, and have to be resolved using a dedicated regulator. After the later regularization, TMDPDFs and TMDWFs acquire an additional rapidity scale dependence. This dependence should cancel in theoretical predictions for physical observables. 

The CS kernel $K\left(b_{\perp}, \mu\right)$, known as the rapidity anomalous dimension, encodes the rapidity dependence  of the TMD distributions~\cite{Collins:1981va,Collins:1981uk}:
\begin{align}
	2 \zeta \frac{d}{d \zeta} \ln f^{\mathrm{TMD}}\left(x, b_{\perp}, \mu, \zeta\right)=K\left(b_{\perp}, \mu\right), \label{eq:TMDsrapidityevolution}
\end{align}
where $f^{\mathrm{TMD}}$ denotes any leading twist TMDPDF or TMDWF. The TMD distributions depend on the longitudinal momentum fraction $x$, transverse spatial separation $b_{\perp}$, which is the Fourier-conjugate to the transverse momentum $k_{\perp}$, as well as the renormalization scale $\mu$ and rapidity scale $\zeta$ {which is related to the hadron momentum}. The $\mu$-dependence of CS kernel $K\left(b_{\perp}, \mu\right)$ satisfies the   renormalization group equation (RGE):
\begin{align}
	\mu^{2} \frac{d}{d \mu^{2}} K\left(b_{\perp}, \mu\right)=-\Gamma_{\text {cusp}}\left(\alpha_{s}\right).  \label{eq:cuspanomalousdimension}
\end{align}
Here $\Gamma_{\text {cusp}}\left(\alpha_{s}\right)=\alpha_sC_F/\pi+\mathcal{O}(\alpha_s^2)$ is the cusp anomalous dimension, which has been calculated  in  perturbation theory  up to 2-loop in Ref.~\cite{Li:2016ctv}, and 3-loop in Ref.~\cite{Moch:2017uml}. The solution to  the RGE  can be expressed as: 
\begin{align}
	K\left(b_{\perp}, \mu\right)=-2 \int_{1 / b_{\perp}}^{\mu} \frac{\mathrm{d} \mu^{\prime}}{\mu^{\prime}} \Gamma_{\text {cusp}}\left(\alpha_{s}\left(\mu^{\prime}\right)\right)+K\left(\alpha_{s}\left(1 / b_{\perp}\right)\right). 
	\label{eq:nonper-per-CS-kernel}
\end{align} 
For large $b_{\perp}$ with $b_{\perp}^{-1}\lesssim\Lambda_{\mathrm{QCD}}$, the CS kernel becomes nonperturbative, which is  represented by the non-cusp anomalous dimension $K\left(\alpha_{s}\left(1/b_{\perp}\right)\right)$ in the above equation. 

In the past decades, the CS kernel has  been widely studied in global fits of  TMD parton distributions~\cite{Landry:1999an,Landry:2002ix,DAlesio:2014mrz,Sun:2014dqm,Konychev:2005iy,Bacchetta:2017gcc,Scimemi:2017etj,Scimemi:2019cmh,Bacchetta:2019sam}. The explicit form in the nonperturbative region can only be parametrized by extending the perturbative expressions at  small $b_{\perp}$, which inevitably introduces systematic uncertainties. A direct calculation of TMDPDFs and the relevant CS kernel 
on the lattice was an almost insurmountable hurdle until the establishment  of LaMET~\cite{Ji:2013dva,Ji:2014gla}.  A remarkable recent development in LaMET is that  these quantities  can  be accessed through the corresponding quasi observables~\cite{Ebert:2019okf,Ji:2019sxk,Ji:2019ewn,Ji:2021znw}.

%%%%%%%%%%%%%%%%%%%%%%%%%%%%%%%


\subsection{Quasi TMD Wave Functions}
As stated above, one can define the quasi TMDWFs for a highly-boosted pseudoscalar meson along the $z$-direction with large momentum $P^z$ as: 
\begin{align}
	\tilde{\Psi}^{\pm}&\left(x,b_{\perp},\mu,\zeta_z\right)=\lim_{L\to\infty}\int\frac{dz}{2\pi}e^{ix_rzP^z}\frac{\tilde{\Phi}^{\pm0}\left(z,b_{\perp},P^z,a,L\right)}{\sqrt{Z_E(2L,b_{\perp},\mu, a)}},\label{eq:quasiWFinmomentumspace}
\end{align}
where $x_r=x-\frac{1}{2}$. The unsubtracted quasi TMDWF $\tilde{\Phi}^{\pm0}$ is defined as an equal-time correlator containing nonlocal quark bilinear operator with staple-shaped gauge link:
\begin{align}
	\tilde{\Phi}^{\pm0}&\left(z,b_{\perp},P^z,a,L\right) =\left\langle 0\right|\bar{\psi}\left(z \hat{n}_{z}/2+b_{\perp} \hat{n}_{\perp}\right) \Gamma \nonumber\\
	&\times U_{\sqsupset, \pm L}\left(z \hat{n}_{z}/2+b_{\perp} \hat{n}_{\perp},-z \hat{n}_{z}/2\right) \psi\left(-z \hat{n}_{z}/2\right)\left| P^{z}\right\rangle. \label{eq:quasiWFincoordinatespace}
\end{align}
For a pseudoscalar mesonic state, the Dirac structure $\Gamma$ can be chosen as $\gamma^z\gamma_5$ or $\gamma^t\gamma_5$, which approaches the leading-twist structure $\gamma^+\gamma_5$ in the light-cone limit. With a large but finite $P^z$, the difference between $\gamma^z\gamma_5$ and $\gamma^t\gamma_5$  is  suppressed by powers of $1/P^z$. Technically, one can also use a combination of them, such as $\left(\gamma^z\gamma_5+\gamma^t\gamma_5\right)/2$ to minimize  power corrections, and more details can be found  in Sec.~\ref{sec:operator}.  Various combinations were also explored in Ref.~\cite{Li:2021wvl}.  The superscript ''0'' in $\tilde{\Phi}^{\pm0}$ indicates bare quantities. The linear divergences come from the self-energy of the gauge-link, 
\begin{align}
	&U_{\sqsupset, \pm L}\left(z \hat{n}_{z}/2+b_{\perp} \hat{n}_{\perp},-z \hat{n}_{z}/2\right) \equiv\nonumber\\
	&	U_z^{\dagger}\left(z \hat{n}_{z}/2+b_{\perp} \hat{n}_{\perp}; L \right) U_{\perp}\big((L-z /2)\hat{n}_{z};{b}_T \big) U_z\left( -z \hat{n}_{z}/2; L \right), \label{eq:stapleshapedUlink}
\end{align}
and does not appear as a pole at $d=4$ in dimensional regularization. The Euclidean gauge link in $U_{\sqsupset, \pm L}$ is defined as
\begin{align}
	U_z(\xi,\pm L)=\mathcal{P}	\exp \left[-i g \int_{\xi^{z}}^{\pm L} d \lambda n_{z} \cdot A\left(\vec{\xi}_{\perp}+n_{z} \lambda\right)\right],
\end{align}
where $\xi^{z}=-\xi \cdot n_{z}$. The $\pm L$  corresponds to the farthest position that the gauge link can reach in positive or negative $n_z$ direction  on a finite Euclidean lattice. This is depicted as the blue and red lines in Fig.~\ref{fig:definitionofquasiTMDWF}.


%%%%%%%%%%%%%%%%%%%%%%%%%%%%%%%%%%
\begin{figure}
\centering
\includegraphics[width=0.5\textwidth]{quasi-TMDWF}
\caption{Illustration of the staple-shaped gauge-link included in unsubtracted quasi-TMDWFs  and related Wilson loop. The blue and red double lines in the upper panel represent the $L$-shift direction on the Euclidean lattice, and the lower panel shows the correspondingly Wilson loop, which will subtract   UV logarithmic  and linear divergences in  quasi-TMDWFs . }
\label{fig:definitionofquasiTMDWF}
\end{figure}
%%%%%%%%%%%%%%%%%%%%%%%%%%%%%%%%%%

Since the  linear divergence is associated with the gauge-link, it can be removed by a similar gauge-link with the same total length. An  optional choice  is to make use the Wilson loop,  denoted as $Z_E$. The Wilson loop can be chosen  as the vacuum expectation of a flat rectangular Euclidean Wilson-loop in the $z$-$\perp$ plane:
\begin{align}\label{eq:WilsonLoop}
Z_E\left(2L, b_{\perp}, \mu,a\right)=\frac{1}{N_{c}} \operatorname{Tr}\left\langle 0\left|U_{\perp}(0;b_{\perp}) U_{z}\left(b_{\perp}\hat{n}_{\perp}; 2L\right)\right| 0\right\rangle. 
\end{align}
Here  the length of $Z_E$ is twice that of the staple-shaped gauge-link $U_z^{(\dagger)}$ in the $z$ direction, and thus it is anticipated that the square root of $Z_E\left(2L, b_{\perp}, \mu\right)$ cancels the linear divergence and heavy quark potential in the gauge-link.  There are residual logarithmic divergences from the vertex of   Wilson line and light quark, which can be renormalized in dimensional regularization~\cite{Ji:2021uvr}. As these logarithmic divergences are independent of $z$, $b_{\perp}$, $P_{z}$ and $L$, they will explicitly cancel out when  the ratio of quasi TMDWFs is studied. 



%%%%%%%%%%%%%%%%%%%%%%%%%%%%%%%%%%%%%%%%
\subsection{Factorization of Quasi TMDWFs}\label{sec:factorizationformula}

With the help of soft function, the infrared contributions in the subtracted quasi TMDWFs can be properly accounted for such that the infrared structures for the   quasi TMDWFs and light-cone ones are matched. This implies a multiplicative factorization theorem in the framework of LaMET~\cite{Ebert:2019okf,Ji:2019sxk,Ji:2019ewn,Ji:2021znw,Ebert:2022fmh}:
\begin{align}
&\tilde \Psi^{\pm}(x, b_{\perp}, \mu, \zeta_z) S_r^{1/2}(b_{\perp}, \mu) \nonumber\\
&= H^{\pm}\left(\zeta_z, \overline\zeta_z, \mu^2\right) \exp\left[\frac{1}{2}K(b_{\perp},\mu)\ln\frac{\mp\zeta_z- i\epsilon}{\zeta} \right]  \nonumber\\
& \qquad\times \Psi^{\pm}(x, b_\perp, \mu, \zeta)+\mathcal{O}\left(\frac{\Lambda_{QCD}^2}{\zeta_z}, M^2\zeta_z,\frac{1}{b_{\perp}^2\zeta_z}\right),
\label{eq:matching}
\end{align}
where the superscript $\pm$ in Eq.~\eqref{eq:matching} corresponds to the direction in the Wilson line, $\Psi^{\pm}$ is the TMDWFs defined in the infinite momentum frame. The reduced soft function $S^{1/2}_r(b_{\perp}, \mu)$ emerges from the different soft gluon radiation effects in $\tilde \Psi^{\pm}$ and $\Psi^{\pm}$~\cite{Ji:2021znw}. The mismatch  of the rapidity scale $\zeta$ and $\zeta_z$ can be compensated by the CS kernel $K(b_{\perp},\mu)$. Both $S$ and $K$ are independent of the $\pm$ choice.  $H^{\pm}$ is the 1-loop perturbative matching kernel~\cite{Ji:2021znw}:
\begin{align}
& H^{\pm}(\zeta_z, \overline{\zeta}_z,\mu) \nonumber\\
& = 1+  \frac{ \alpha_s C_F}{4\pi} \bigg(-\frac{5\pi^2}{6}-4 +{\ell}_{\pm} + \overline \ell_{\pm} - \frac{1}{2} ({\ell}_{\pm}^2 + \overline \ell_{\pm}^2) \bigg).  \label{eq:1loophardkernel}
\end{align}
With the abbreviations ${\ell}_{\pm} = \ln \left[(-\zeta_z \pm i\epsilon)/\mu^2\right]$, and $\overline {\ell}_{\pm} = \ln \left[(-\overline \zeta_z \pm i\epsilon)/\mu^2\right]$, the scales $\zeta_z = (2x P^z)^2$ and $\overline\zeta_z = \left(2\bar{x} P^z\right)^2$, and $\bar{x}=1-x$. It should be noticed that $H^{\pm}$ contains nonzero imaginary parts in $\ell_{\pm}$ and $\bar{\ell}_{\pm}$. While the imaginary parts in  $\ell_{\pm}/\bar{\ell}_{\pm}$ are constants, the  ones in the double logarithms ${\ell}_{\pm}^2/\bar{\ell}_{\pm}^2$ are momentum-dependent.

A characteristic behavior of Eq.(\ref{eq:matching}) is that this factorization is multiplicative~\cite{Ji:2020ect}, which indicates that hard gluon contributions are local. This is due to the fact that hard gluon exchange between  the quark and anti-quark  sectors in quasi TMDWFs is power suppressed: if there were such a hard gluon, the spatial separation between its attachments  is much smaller than $b_{\perp}$, resulting in power suppression compared to the typical hard mode contributions.  
Thus, at leading power the factorization of quasi TMDWFs are multiplicative. This feature  is illustrated in Fig.\ref{fig:reduced_graph}, in which the collinear, soft and hard sub-diagrams represent the pertinent  contributions. Further, $\zeta_z$ arising from   Lorentz-invariant combinations of collinear momentum modes, will provide the natural hard scale of the hard sub-diagram. More detailed explanations for the factorization of quasi TMDPDFs in LaMET are given in the recent review~ \cite{Ji:2019ewn}. 


%%%%%%%%%%%%%%%%%%%%%%%%%%%%%%%%%%
\begin{figure}
\centering
\includegraphics[width=0.5\textwidth]{reduced_graph}
\caption{Leading-power reduced graph for pseudoscalar meson quasi TMDWFs. Here $C, ~S$ denote the collinear and soft sectors of the infra-red structure, while $H$ denote the hard contributions. Since the hard-gluon exchange between the quark and anti-quark is power suppressed, the hard parts are disconnected with each others, and therefore the factorization of quasi-WF amplitude is multiplicative.  }
\label{fig:reduced_graph}
\end{figure}
%%%%%%%%%%%%%%%%%%%%%%%%%%%%%%%%%%




\subsection{Collins-Soper Kernel From Quasi TMDWFs}\label{sec:extractingCSkernel}


From Eq.(\ref{eq:matching}), one can see that the momentum dependence in   quasi TMDWFs provides an option to determine the CS kernel. This can be written  in a way similar to Eq.(\ref{eq:TMDsrapidityevolution})~\cite{Ji:2021znw}: 
\begin{align}
	2 \zeta_z \frac{d}{d \zeta_z} \ln \tilde \Psi^{\pm}\left(x, b_{\perp}, \mu, \zeta_z\right)=&K\left(b_{\perp}, \mu\right) +\frac{1}{2}\mathcal{G}^{\pm}\left(x^2\zeta_z,\mu\right) \nonumber\\
	+& \frac{1}{2}\mathcal{G}^{\pm}\left(\bar{x}^2\zeta_z,\mu\right)+\mathcal{O}\left(\frac{1}{\zeta_z}\right), \label{eq:quasiTMDsrapidityevolution}
\end{align}
where $K\left(b_{\perp}, \mu\right)$ denotes the same kernel as in Eq.(\ref{eq:TMDsrapidityevolution}), and does not depend on the hard scale $\zeta_z$ for large $P^z$. Unlike TMDWFs, the  quasi distributions also contain hard contributions, whose  rapidity dependence is  represented by the perturbative  $\mathcal{G}^{\pm}$ as a function of hard scale $\zeta_z$. 
From the $\zeta_z$ dependence of quasi TMDWFs, we can see that when $P^z\to\infty$, the large logarithms in $P^z$ are partially absorbed into $K\left(b_{\perp}, \mu\right)$  and the remanent is incorporated in the perturbative matching kernel. Therefore,  both the matching kernel $H^{\pm}$  and an exponential of  CS kernel $K\left(b_{\perp}, \mu\right)$ are needed to describe the dependence on $\zeta_z$ of quasi TMDWFs. 

In order to extract the CS kernel $K\left(b_{\perp}, \mu\right)$ explicitly, one  can make use of    Eq.(\ref{eq:matching}) with  two different large momenta $P_1^z\neq P_2^z\gg1/b_{\perp}$ but the same   scale $\zeta_z$. Taking a  ratio of these two quantities  gives 
\begin{align}
	\frac{\tilde \Psi^{\pm}(x, b_{\perp}, \mu,P_1^z)}{\tilde \Psi^{\pm}(x, b_{\perp}, \mu, P_2^z)}=\frac{H^{\pm}\left(xP_1^z,\mu \right)}{H^{\pm}\left(xP_2^z,\mu \right)}\exp\left[K(b_{\perp},\mu) \ln\frac{P_1^z}{P_2^z}\right], \label{eq:ratioofquasiWF}
\end{align}
where the reduced soft function $S_r(b_{\perp},\mu)$ and TMDWFs $\Psi^{\pm}(x,b_{\perp},\mu,\zeta)$  have been canceled in the ratio. Therefore, the CS kernel $K\left(b_{\perp}, \mu\right)$ can be extracted through 
\begin{align}
	K\left(b_{\perp}, \mu\right)=\frac{1}{\ln(P_1^z/P_2^z)} \ln\frac{H^{\pm}(xP_2^z,\mu)\tilde{\Psi}^{\pm}(x,b_{\perp},\mu,P_1^z)}{H^{\pm}(xP_1^z,\mu)\tilde{\Psi}^{\pm}(x,b_{\perp},\mu,P_2^z)}. \label{eq:extractingCSkernel}
\end{align}
Note that the extracted result is formally independent of $x$ and $P_{1/2}^z$ at leading power,  and both $\tilde{\Psi}^{+}$ and $\tilde{\Psi}^{-}$ can be used to extract $K\left(b_{\perp}, \mu\right)$.  This is derived at the leading power in the factorization scheme and  might be undermined by power corrections. Accordingly, in order  to reduce  the systematic uncertainties,  we take the average: 
 \begin{align}
	K\left(b_{\perp}, \mu\right)=&\frac{1}{2\ln(P_1^z/P_2^z)} \left[\ln\frac{H^{+}(xP_2^z,\mu)\tilde{\Psi}^{+}(x,b_{\perp},\mu,P_1^z)}{H^{+}(xP_1^z,\mu)\tilde{\Psi}^{+}(x,b_{\perp},\mu,P_2^z)}\right.  \nonumber\\
	&+\left.\ln\frac{H^{-}(xP_2^z,\mu)\tilde{\Psi}^{-}(x,b_{\perp},\mu,P_1^z)}{H^{-}(xP_1^z,\mu)\tilde{\Psi}^{-}(x,b_{\perp},\mu,P_2^z)}\right].
\end{align}
The details will be discussed in Sec. \ref{subsec:csresults}.
